\documentclass[aspectratio=169]{beamer}

\usepackage[T1]{fontenc}
\usepackage[utf8]{inputenc}
\usepackage{booktabs}
\usepackage{array}
\usepackage{tabularx}
\usepackage{graphicx}
\usepackage{xcolor}

\definecolor{evblue}{HTML}{1F4E79}
\definecolor{evgreen}{HTML}{2E7D32}
\definecolor{evorange}{HTML}{C55A11}
\definecolor{evgray}{HTML}{5F6B73}
\definecolor{evlight}{HTML}{F3F6F8}

\setbeamertemplate{navigation symbols}{}
\setbeamercolor{title}{fg=evblue}
\setbeamercolor{frametitle}{fg=evblue}
\setbeamercolor{structure}{fg=evblue}
\setbeamercolor{block title}{fg=white,bg=evblue}
\setbeamercolor{block body}{fg=black,bg=evlight}
\setbeamerfont{frametitle}{series=\bfseries}

\newcolumntype{Y}{>{\raggedright\arraybackslash}X}

\newcommand{\kw}{\mathrm{kW}}
\newcommand{\mw}{\mathrm{MW}}
\newcommand{\kwh}{\mathrm{kWh}}
\newcommand{\mwh}{\mathrm{MWh}}
\newcommand{\gwh}{\mathrm{GWh}}
\newcommand{\takeaway}[1]{\vspace{0.5em}\textcolor{evblue}{\textbf{Takeaway:}} #1}
\newcommand{\caveat}[1]{\textcolor{evorange}{\textbf{Caveat:}} #1}
\newcommand{\maybegraphic}[2][]{%
  \IfFileExists{#2}{%
    \includegraphics[#1]{#2}%
  }{%
    \setlength{\fboxsep}{1em}%
    \fbox{%
      \begin{minipage}[c][0.42\textheight][c]{0.82\linewidth}%
        \centering
        \textcolor{evgray}{\textbf{Figure not available}}\\[0.4em]
        \footnotesize\texttt{\detokenize{#2}}%
      \end{minipage}%
    }%
  }%
}

\title[Albany EV Charging Results]{Albany County Behavioral EV Charging Results}
\subtitle{Full dataset run with ZIP/ZCTA adoption scaling, H3-8 charging geography, and GridUp-style charger availability}
\author{NYSERDA EV Adoption and Charging Model}
\date{June 23, 2026}

\begin{document}

% -----------------------------------------------------------------------------
\begin{frame}
  \titlepage
\end{frame}

\begin{frame}{This week's status}
\small
\begin{tabularx}{\linewidth}{p{0.28\linewidth}Y}
\toprule
\textbf{Area} & \textbf{Current status} \\
\midrule
Model geography & Adoption now scales from home ZIP/ZCTA; charging load is placed at destination H3. \\
Charging behavior & Availability now follows the GridUp-style inverse-demand choice model; allocation still uses conditional expected value across ranked stops. \\
Managed charging & Applied only to managed-eligible home/work L2 sessions; public, civic, and DCFC sessions remain unmanaged. \\
Full run & Full Albany dataset completed with H3 required and ZIP/ZCTA adoption scaling enabled. \\
Visualization & Albany ZIP adoption, H3 charging-demand, and selected-H3 point maps rebuilt from the current forecast/model outputs. \\
\bottomrule
\end{tabularx}
\end{frame}

% -----------------------------------------------------------------------------
\begin{frame}{Remaining trip-chain modeling decision: final stops}
\small
\textbf{Final activity distribution for chains that do not end at home}

\vspace{0.4em}
\begin{tabular}{lr}
\toprule
\textbf{Final activity type} & \textbf{Chains} \\
\midrule
Social / Recreational & 7,572 \\
Work & 3,300 \\
School / Daycare / Religious & 1,620 \\
Shopping / Errands & 909 \\
Other & 715 \\
Transport / drop-off / pick-up & 450 \\
Meals & 117 \\
Medical / Dental & 105 \\
\bottomrule
\end{tabular}

\begin{itemize}
  \item \textbf{As observed:} final non-home stops remain non-home overnight stops.
  \item \textbf{Completed chain:} add a synthetic return-home leg and overnight home dwell.
\end{itemize}
\end{frame}

% -----------------------------------------------------------------------------
\begin{frame}{Geography design: home adoption vs. physical load}
\small
\begin{block}{Core decision}
Keep home adoption geography separate from charging-load geography.
\end{block}

\begin{tabularx}{\linewidth}{p{0.26\linewidth}p{0.26\linewidth}Y}
\toprule
\textbf{Use case} & \textbf{Target geography} & \textbf{Reason} \\
\midrule
EV adoption forecast & Home ZIP/ZCTA & Adoption model is ZIP-level; join from home coordinate to ZIP/ZCTA. \\
Physical charging load & Destination H3 & Load appears where vehicles charge, not necessarily where they live. \\
\bottomrule
\end{tabularx}
\end{frame}


% -----------------------------------------------------------------------------
\begin{frame}{GridUp-style availability and charging allocation}
\footnotesize
\setlength{\abovedisplayskip}{0.25em}
\setlength{\belowdisplayskip}{0.25em}
\begin{block}{Availability model}
Configured probabilities are interpreted as peak stopped-vehicle-hour availability, then adjusted by local stopped-vehicle demand:
\[
q_{j,t}=\min\left(1,\;q^{\text{peak}}_j
\frac{N^{\text{peak}}_{\ell(j)}}{N_{\ell(j),t}}\right)
\]
where $\ell(j)$ is charger assumption, charger type, and rounded destination point.
\end{block}

\begin{block}{Allocation method}
Carry probability states across ranked charging opportunities:
\[
\mathcal{S}_j = \{(p_k, r_k): \text{probability } p_k, \text{ remaining energy } r_k\}
\]
At each stop, split every state into available and unavailable branches, allocate:
\[
E_{j,k}=\min(r_k, C_j)
\]
and aggregate expected event energy:
\[
\mathbb{E}[E_j]=\sum_k p_k \cdot q_j \cdot E_{j,k}
\]
Stop ranking is home, work, long public, civic, then quick public.
\end{block}
\end{frame}

% -----------------------------------------------------------------------------
\begin{frame}{Managed vs. unmanaged charging}
\footnotesize
\begin{columns}[T,totalwidth=\linewidth]
\begin{column}{0.48\linewidth}
\begin{block}{Unmanaged}
Charging begins at arrival and uses charger power until allocated energy is delivered or the dwell window closes:
\[
T_s = \min\left(\tau_s, \frac{E_s}{P_s}\right)
\]
\end{block}
\end{column}
\begin{column}{0.48\linewidth}
\begin{block}{Managed}
Only managed-eligible charger types are spread across the dwell window:
\[
\bar{P}_s = \min\left(P_s, \frac{E_s}{\tau_s}\right)
\]
\end{block}
\end{column}
\end{columns}

\vspace{0.3em}
\begin{itemize}
  \item Peak-hour availability assumptions: single-family home L2 1.00, multifamily home L2 0.50, work L2 0.50, civic L2 0.50, long public L2 1.00, quick public DCFC 0.35.
  \item GridUp adjustment raises availability in lower-demand stopped-vehicle hours and clips probability at 1.00.
  \item This keeps home/work management separate from public/DCFC charging behavior.
\end{itemize}
\end{frame}

% -----------------------------------------------------------------------------
\begin{frame}{Adoption scaling logic}
\small
\begin{block}{Forecast scaling equation}
\[
L_{g,t,c,y} = \sum_h L^{\text{fullEV}}_{h,g,t,c}
\cdot A_{h,y,s} \cdot G_{h,y,s}
\]
\end{block}

\begin{tabularx}{\linewidth}{p{0.16\linewidth}Y}
\toprule
\textbf{Symbol} & \textbf{Meaning} \\
\midrule
$h$ & Home ZIP/ZCTA used for adoption attribution. \\
$g$ & Charging H3 cell where load physically appears. \\
$c$ & Charger category / charging location type. \\
$A_{h,y,s}$ & Adoption fraction for home geography, year, and scenario. \\
$G_{h,y,s}$ & Vehicle growth factor for home geography, year, and scenario. \\
\bottomrule
\end{tabularx}
\end{frame}

% -----------------------------------------------------------------------------
\begin{frame}{Result Summary}
\small
\begin{tabular}{p{0.37\linewidth}p{0.52\linewidth}}
\toprule
\textbf{Question} & \textbf{Current result} \\
\midrule
Full-EV daily energy delivered & 3.723 \(\gwh\) per modeled day under both managed and unmanaged modes. \\
Full-EV peak load & 337.3 \(\mw\) unmanaged at 20:00; 229.8 \(\mw\) managed at 03:00. \\
2027 baseline adoption load & 174.6 \(\mwh\) per modeled day; 15.81 \(\mw\) unmanaged peak. \\
Managed charging effect & About 31.9\% lower full-EV peak and 31.9\% lower 2027 baseline peak. \\
\bottomrule
\end{tabular}

\takeaway{The model now produces a home-ZCTA by charging-H3 by time load surface, then scales it with the ZIP-level adoption forecast.}
\end{frame}

% -----------------------------------------------------------------------------
\begin{frame}{Modeling Geography}
\small
\begin{columns}[T,totalwidth=\linewidth]
\begin{column}{0.49\linewidth}
\textbf{Adoption scaling geography}
\begin{itemize}
  \item Adoption forecast is ZIP/ZCTA-level.
  \item Each modeled vehicle day keeps its home ZCTA.
  \item Scaling is applied by \texttt{home\_zcta}, forecast year, and adoption scenario.
\end{itemize}
\end{column}
\begin{column}{0.49\linewidth}
\textbf{Charging-load geography}
\begin{itemize}
  \item Charging load is placed at the destination charging location.
  \item Primary planning geography is H3 at resolution 8.
\end{itemize}
\end{column}
\end{columns}
\end{frame}

% % -----------------------------------------------------------------------------
% \begin{frame}{Output Tables Produced}
% \scriptsize
% \renewcommand{\arraystretch}{1.15}
% \begin{tabular}{p{0.72\linewidth}r}
% \toprule
% \textbf{What the table contains} & \textbf{Rows} \\
% \midrule
% Expected charging sessions: timing, energy, charger type, event probability, and home/destination geography. & 678,768 \\
% Vehicle-day energy accounting by managed/unmanaged mode: daily energy need, delivered energy, and unmet energy. & 433,920 \\
% Pre-adoption load surface by home ZCTA, charging H3, charger type, and 15-minute time bin. & 1,230,443 \\
% Hourly version of the pre-adoption home-ZCTA by charging-H3 load surface. & 380,541 \\
% Full-EV charging-H3 load rollups used for spatial peak and profile analysis. & 502,469 / 140,025 \\
% Destination-point load rollups retained for QA and possible future feeder or transformer joins. & 7,410,594 / 2,037,569 \\
% Adoption-scaled charging-H3 planning loads across forecast years and adoption scenarios. & 15,074,070 / 4,200,750 \\
% \bottomrule
% \end{tabular}
% \end{frame}

% -----------------------------------------------------------------------------
\begin{frame}{Full-EV Energy Accounting}
\small
\begin{tabular}{p{0.43\linewidth}r}
\toprule
\textbf{Metric} & \textbf{Value per modeled day} \\
\midrule
Vehicle days & 216,960 \\
Total network miles & 8.964 million miles \\
Mean daily network miles & 41.3 miles per vehicle day \\
Daily travel energy need & 3.765 \(\gwh\) \\
Expected charging energy delivered & 3.723 \(\gwh\) \\
Unmet energy & 0.042 \(\gwh\) \\
Unmet share of daily travel energy & 1.12\% \\
\bottomrule
\end{tabular}

\vspace{0.8em}
\begin{block}{Interpretation}
The unmet-energy metric is not a grid load. It is a behavioral feasibility signal showing where the current stop, dwell, charger-access, and power assumptions cannot fully replenish modeled daily travel energy.
\end{block}
\end{frame}

% -----------------------------------------------------------------------------
\begin{frame}{Full-EV Hourly Load Shape}
\begin{center}
\maybegraphic[width=0.86\linewidth]{figures/full_ev_hourly_profile.pdf}
\end{center}

\small
\takeaway{Managed charging shifts the aggregate peak from 20:00 to 03:00 and reduces the full-EV peak from 337.3 \(\mw\) to 229.8 \(\mw\).}
\end{frame}

% -----------------------------------------------------------------------------
\begin{frame}{Charging Energy Is Concentrated at Home}
\begin{center}
\maybegraphic[width=0.74\linewidth]{figures/energy_share_2027_baseline.pdf}
\end{center}

\small
\takeaway{Home charging accounts for 87.4\% of 2027 baseline daily energy; long public contributes 7.1\%, work 3.8\%, quick public 1.3\%, and civic 0.5\%.}
\end{frame}

% -----------------------------------------------------------------------------
\begin{frame}{2027 Baseline Adoption-Scaled Load}
\begin{center}
\maybegraphic[width=0.86\linewidth]{figures/baseline_2027_hourly_profile.pdf}
\end{center}

\small
\takeaway{For 2027 baseline adoption, managed charging lowers the county-wide hourly peak from 15.81 \(\mw\) to 10.76 \(\mw\), with the same total daily energy of 174.6 \(\mwh\).}
\end{frame}

% -----------------------------------------------------------------------------
\begin{frame}{Adoption-Scaled Energy by Forecast Year}
\begin{center}
\maybegraphic[width=0.86\linewidth]{figures/scaled_energy_by_year.pdf}
\end{center}
\end{frame}

% -----------------------------------------------------------------------------
\begin{frame}{Adoption-Scaled Peak by Forecast Year}
\begin{center}
\maybegraphic[width=0.86\linewidth]{figures/scaled_peak_by_year.pdf}
\end{center}
\end{frame}

% -----------------------------------------------------------------------------
\begin{frame}{Interactive Visualization Outputs}
\small
\renewcommand{\arraystretch}{1.15}
\begin{tabularx}{\linewidth}{p{0.31\linewidth}Y}
\toprule
\textbf{Visualization} & \textbf{Purpose} \\
\midrule
ZIP adoption forecast map & Shows forecast EV stock or adoption share by home ZIP/ZCTA, scenario, and date. This is the adoption side of the workflow. \\
H3 charging-demand map & Shows adoption-scaled charging demand by destination H3 cell, scenario, forecast year, managed mode, metric, and hour. This is the physical-load side of the workflow. \\
Selected-H3 point map & Drills into one H3 cell and shows the destination points that create load inside that cell. This supports QA and future feeder or transformer joins. \\
\bottomrule
\end{tabularx}

\vspace{0.7em}
\scriptsize
\textbf{Files:} \texttt{visualization/albany\_zip\_adoption\_forecast\_map.html}, \texttt{visualization/albany\_h3\_charging\_demand\_forecast\_map.html}, and selected-H3 point map/CSV outputs.
\end{frame}

% -----------------------------------------------------------------------------
\begin{frame}{Visualization: Adoption Forecast}
\begin{center}
\maybegraphic[height=0.60\textheight]{figures/albany_zip_adoption_forecast_map.png}
\end{center}

\small
\takeaway{This view is for checking the home-side adoption forecast before it is used to scale the behavioral charging surface.}
\end{frame}

% -----------------------------------------------------------------------------
\begin{frame}{Visualization: H3 Charging Demand Forecast}
\begin{center}
\maybegraphic[height=0.60\textheight]{figures/albany_h3_charging_demand_forecast_map.png}
\end{center}

\small
\takeaway{This view converts the modeled behavior into spatial grid demand: scenario, year, managed mode, metric, and hour can be inspected at H3-8 resolution.}
\end{frame}

% -----------------------------------------------------------------------------
\begin{frame}{Visualization: Point-Level Demand Within One H3 Cell}
\begin{center}
\maybegraphic[height=0.60\textheight]{figures/albany_h3_point_charging_demand_map.png}
\end{center}

\small
\takeaway{The point-level view explains which destination points are driving load inside a selected H3 cell, which is useful before linking demand to utility assets.}
\end{frame}

% -----------------------------------------------------------------------------
\begin{frame}{2027 Baseline H3 Hotspots}
\small
\begin{columns}[T,totalwidth=\linewidth]
\begin{column}{0.49\linewidth}
\textbf{Unmanaged top H3-hours}

\vspace{0.3em}
\scriptsize
\begin{tabular}{lrr}
\toprule
\textbf{H3 cell} & \textbf{Hour} & \textbf{\(\kw\)} \\
\midrule
882b8902b1fffff & 20 & 165.6 \\
882b8902b1fffff & 19 & 163.2 \\
882b890025fffff & 19 & 160.6 \\
882b8902b1fffff & 18 & 159.9 \\
882b890025fffff & 18 & 156.4 \\
\bottomrule
\end{tabular}
\end{column}
\begin{column}{0.49\linewidth}
\textbf{Managed top H3-hours}

\vspace{0.3em}
\scriptsize
\begin{tabular}{lrr}
\toprule
\textbf{H3 cell} & \textbf{Hour} & \textbf{\(\kw\)} \\
\midrule
882b8902b1fffff & 3 & 107.4 \\
882b8902b1fffff & 1 & 107.1 \\
882b8902b1fffff & 0 & 106.6 \\
882b8902b1fffff & 4 & 106.6 \\
882b8902b1fffff & 2 & 106.6 \\
\bottomrule
\end{tabular}
\end{column}
\end{columns}
\end{frame}

% -----------------------------------------------------------------------------
\begin{frame}{Current Caveats}
\small
\begin{itemize}
  \item Albany trip data supports destination coordinates, purpose, and dwell time, but not detailed destination land use.
  \item Office vs. non-office workplace charging and DCFC 150/250/350 kW splits are still approximated with configurable assumptions.
  \item Charger-access and charger-power assumptions are configurable but not yet calibrated against local charger inventory and housing data.
  \item The adoption forecast in this run covers 2018--2027, so longer-term grid planning requires a longer adoption horizon.
\end{itemize}
\end{frame}

% -----------------------------------------------------------------------------
\begin{frame}{Next Steps}
\small
\begin{enumerate}
  \item Map charging H3 cells to feeders, substations, or utility service polygons when those layers become available.
  \item Join destinations to land-use and EVSE inventory layers to calibrate workplace and DCFC charger choices.
  \item Calibrate synthetic vehicle-day weights to observed Albany County vehicle stock and EV registrations.
  \item Extend the ZIP/ZCTA adoption forecast and expand daily profiles to full-year load shapes for planning studies.
  \item Run sensitivity cases for charger access, managed eligibility, DCFC mix, and energy efficiency.
\end{enumerate}
\end{frame}

\end{document}
